% Chapter 8: Limitations and Future Work

\chapter{Limitations and Future Work}

This chapter expands on the limitations of the current PVE-ZKP prototype and
outlines directions for future research and engineering work. The goal is to
clarify what the present system does \emph{not} provide and to sketch a roadmap
for strengthening and extending it.

\section{Trusted Setup}

The prototype relies on a single-party trusted setup for each Groth16 circuit.
This is acceptable for experimentation but undesirable in production, as a
setup participant who retains trapdoor information could, in principle,
construct fraudulent proofs.

Future work should replace the single-party setup with a multi-party
computation (MPC) ceremony, ideally using existing tooling and best practices
from other SNARK deployments. The monograph does not prescribe a specific
ceremony design but treats this as a modular component.

\section{Proof Storage and Durability}

Proofs are currently stored in an in-memory map inside the validation service.
This design is simple but has several drawbacks: proofs are lost on service
restart, no explicit retention policy is enforced, and there is no access
control beyond the service's own process boundary.

A more robust design would use a dedicated data store such as a database or
cache with strict access control, audit logging, and time-to-live policies.
This would also support horizontal scaling and fault tolerance for the
validation service.

\section{Performance and Scalability}

The prototype does not include empirical benchmarks for proof-generation time,
verification time, or resource consumption. These measurements are essential
for understanding how PVE-ZKP behaves under realistic loads and for identifying
bottlenecks.

Future work should design and run experiments that measure performance for
various circuit sizes, concurrency levels, and deployment settings. Results
would inform circuit design, parameter choices, and the feasibility of
large-scale deployments.

\section{Side-Channel Leakage}

Timing differences, error messages, and network-level patterns may reveal
information about inputs beyond what is captured in the zero-knowledge proofs
themselves. The prototype does not attempt to equalize computation time or
randomize responses to mitigate such side channels.

Investigating the impact of side-channel leakage in this architecture, and
exploring mitigations such as constant-time implementations or rate limiting,
is an important direction for future work.

\section{Policy Coverage and Expressiveness}

The five policy circuits implemented in PVE-ZKP cover only a subset of
validation patterns that arise in real applications. They do not, for example,
handle complex cross-field dependencies, multi-step workflows, or adaptive
policies based on external data.

Extending the policy library to support richer constraints, while keeping
circuits auditable and efficient, is an open design challenge. Additional work
is needed to develop higher-level abstractions for composing policy circuits
and for verifying that combined policies behave as intended.

\section{Toward a Full System and Thesis}

Finally, this monograph represents an intermediate step between a prototype and
a full thesis. Future work includes not only the technical enhancements
outlined above but also more formal modeling, security proofs where appropriate,
and comprehensive evaluation across multiple application domains.
